\section{\ActivationCircuits{}}

We now study schemes to represent the activation network, which enable us to derive sampling schemes from \CompActNets{}.
To this end, we augment the activation network with ancilla variables, understanding it as a conditioned Bayesian network.
We then apply controlled rotations to find q-samples.

\subsection{Uniformly controlled rotations}

\ActivationCircuits{} are uniformly controlled unitaries (see \cite{mottonen_transformation_2005}), where the rotation axis is chosen as the $y$-axes of the Bloch sphere, and the angles are computed by the function $h(\cdot)$ acting on a function value to be encoded.

%% Angle preparing
We define the angle preparing function on $p\in[0,1]$ by
\begin{align*}
    h(p) = 2 \cdot \mathrm{cos}^{-1}\left(\sqrt{1-p}\right) \, .
\end{align*}
For any $p\in[0,1]$ we then have
\begin{align*}
    \contractionof{\onehotmapofat{0}{\cavariable{\insymbol}},\yrotationofat{h(p)}{\cavariable{\insymbol},\cavariable{\outsymbol}}}{\cavariable{\outsymbol}}
    = \coloredmatrixof{\sqrt{1-p} \\ \sqrt{p}}{\cavariable{\outsymbol}} \, .
\end{align*}

%\begin{definition}[\ActivationCircuit{}]
%    Given a function
%    \begin{align*}
%        \exformula : \facstates \rightarrow [0,1]
%    \end{align*}
%    its \activationCircuit{} is the uniformly controlled unitary
%    $\qcaencodingofat{\exformula}{\cavariable{\insymbol},\cavariable{\outsymbol},\shortcatvariables}$ defined as
%    \begin{align*} %% Could ease the definition by coordinates
%        \qcaencodingofat{\hypercore}{\cavariable{\insymbol},\cavariable{\outsymbol},\shortcatvariables}
%        = \sum_{\shortcatindicesin} \onehotmapofat{\shortcatindices}{\shortcatvariables}
%        \otimes \yrotationofat{h(\exformulaat{\shortcatindices})}{\cavariable{\insymbol},\cavariable{\outsymbol}} \, .
%    \end{align*}
%\end{definition}

%% Probability Tensor
When we have a probability tensor, its \activationCircuit{} be prepared, since all values are in $[0,1]$.
Note that for rejection sampling, only the quotients of the values are important, we can therefore scale the value by a scalar such that the mode is $1$.

%% Rescaling tensors with larger coordinates
Tensors $\hypercoreat{\shortcatvariables}$ with non-negative coordinates can be encoded after dividing them by their maximum, that is the \activationCircuit{} of the function
\begin{align*}
    \exformulaat{\shortcatindices}
    = \frac{\hypercoreat{\shortcatvariables=\shortcatindices}}{\max_{\secheadindexof{[\seldim]}} \hypercoreat{\shortcatvariables=\secheadindexof{[\seldim]}}}
\end{align*}
When the maximum of the tensor is not known, it can be replaced by an upper bound (reducing the acceptance rate of the rejection sampling).

%%We revisit the correspondence of Bayesian Networks with graph-controlled circuits, by showing that any Bayesian Network can be prepared by \activationCircuits{} of the conditional probability distributions.
%Q-samples of Bayesian Networks can be prepared by quantum circuits, where each conditional probability distribution is prepared by an \activationCircuit{}.
%
%\begin{theorem}[\cite{low_quantum_2014}]
%    Let there be a Bayesian Network $\probwith$ on a directed acyclic hypergraph $\graph=([\catorder],\edges)$ with dimensions $\catdimof{\catenumerator}=2$ and decorations of the edges $\edges=\{(\parentsof{\catenumerator},\{\catenumerator\}) \wcols \catenumeratorin\}$ by conditional probabilities
%    \begin{align*}
%        \condprobof{\catvariableof{\catenumerator}}{\catvariableof{\parentsof{\catenumerator}}} \, .
%    \end{align*}
%    To each $\catenumeratorin$ we construct the function
%    \begin{align*}
%        \exfunctionof{\catenumerator} \defcols \bigtimes_{\seccatenumerator\in\parentsof{\catenumerator}} [2] \rightarrow [2]
%        \quad,\quad
%        \exfunctionofat{\catenumerator}{\catindexof{\parentsof{\catenumerator}}}
%        = \condprobof{\catvariableof{\catenumerator}=1}{\indexedcatvariableof{\parentsof{\catenumerator}}} \, .
%    \end{align*}
%    Then the concatenation of the \activationCircuits{} (see \figref{fig:cpdEncoding})
%    \begin{align*}
%        \qcaencodingofat{
%            \exfunctionof{\catenumerator}
%        }{\ccatvariableof{\catenumerator}{\insymbol},\catvariableof{\catenumerator},\catvariableof{\parentsof{\catenumerator}}}
%    \end{align*}
%    respecting the order $[\catorder]$ prepares the q-sample to the Bayesian Network $\probwith$ when acting on the initial state $\bigotimes_{\atomenumeratorin}\onehotmapofat{0}{\ccatvariableof{\catenumerator}{\insymbol}}$.
%\end{theorem}
%\begin{proof}
%    Let $\qstateat{\shortcatvariables}$ be the state prepared by the \activationCircuits{} to the conditional probability distributions.
%    We need to show that the measurement distribution of $\qstate$ is the Bayesian Network.
%    For any $\shortcatindices\in\atomstates$ we have
%    \begin{align*}
%        \absof{\qstateat{\indexedshortcatvariables}}^2
%        &= \prod_{\catenumeratorin}
%        \left(\qcaencodingofat{
%            \exfunctionof{\catenumerator}
%        }{\ccatvariableof{\catenumerator}{\insymbol}=0,\indexedcatvariableof{\catenumerator},\indexedcatvariableof{\parentsof{\catenumerator}}}\right)^2 \\
%        &= \prod_{\catenumeratorin}
%        \condprobof{\indexedcatvariableof{\catenumerator}}{\indexedcatvariableof{\parentsof{\catenumerator}}} \\
%        & = \probat{\indexedshortcatvariables} \, .
%    \end{align*}
%\end{proof}

%\begin{figure}
%    \begin{center}
%        \input{./tikz_pics/directed_to_circuit.tex}
%    \end{center}
%    \caption{
%        Representation of directed and positive tensor by a controlled rotation.
%        a) Conditional probability tensor $\condprobat{\catvariableof{\catenumerator}}{\catvariableof{\parentsof{\catenumerator}}}$ being a tensor in a Bayesian Network.
%        b) Circuit Encoding as a controlled rotation, which is the \ActivationCircuit{} of the tensor $\condprobat{\catvariableof{\catenumerator}=1}{\catvariableof{\parentsof{\catenumerator}}}$.
%    }\label{fig:cpdEncoding}
%\end{figure}

%\input{./examples/graph-controlled-circuits/grover-rudolph}

\subsection{Ancilla augmentation}

The computation network consists already of directed cores and therefore is already a Bayesian network.
The activation network however needs ancilla augmentation.
Let us first define the ancilla augmentation of a distribution, and then find ancilla augmentations to \CompActNets{}, which are Bayesian networks.

\begin{definition}[Ancilla Augmentation of a Distribution]
    Let $\probwith$ be a probability distribution of variables $\shortcatvariables$.
    Another distribution $\secprobat{\shortcatvariables,\avariables}$ of variables $\shortcatvariables$ and ancilla variables $\avariables$ is called an ancilla augmentation of $\probwith$, if there is $\acceptanceprob>0$ with
    \begin{align*}
        \acceptanceprob \cdot \probwith
        = \secprobat{\shortcatvariables,\avariables=1_{[\seldim]}} \, .
    \end{align*}
    We refer to $\acceptanceprob$ as the corresponding acceptance probability.
\end{definition}

%% Rejection Sampling, interpretation of \acceptanceprob
Sampling from the distribution can be done by rejection sampling, also called post selection, on the ancilla augmented distribution in the following way.
When drawing a sample $(\shortcatindices,\aindex)$ from $\secprobat{\shortcatvariables,\avariable}$ and rejecting it whenever $\aindex=0$, the accepted samples are distributed as $\probwith$.
We then have $\acceptanceprob$ as the probability of accepting a sampling in this rejection sampling method.
The acceptance probability is critical in determining the efficiency of this approach, since we need to draw $\sim\frac{1}{\acceptanceprob}$ samples from $\secprobat{\shortcatvariables,\avariable}$ to get an accepted sample.

%%% Augmentation cores
%We perform the ancilla augmentation on \ComputationActivationNetworks{} by replacing each activation core $\acttensorofat{\edge}{\headvariableof{\edge}}$ with the augmented core
%\begin{align*}
%    \anaugmentationofat{\edge}{\avariableof{\edge},\headvariableof{\edge}}
%    = \sum_{\catindexof{\edge}} \onehotmapofat{\catindexof{\edge}}{\headvariableof{\edge}} \otimes \Big(
%    \frac{\acttensorofat{\edge}{\indexedheadvariableof{\edge}}}{
%        \max_{\secheadindexof{\edge}}\acttensorofat{\edge}{\indexedheadvariableof{\edge}}
%    } \tbasisat{\avariableof{\edge}}
%    + \big(
%    1 - \frac{\acttensorofat{\edge}{\indexedheadvariableof{\edge}}}{
%        \max_{\secheadindexof{\edge}}\acttensorofat{\edge}{\indexedheadvariableof{\edge}}
%    }
%    \big) \fbasisat{\avariableof{\edge}}
%    \Big)
%\end{align*}


\begin{definition}
    The \activationCircuit{} of a non-negative tensor $\acttensorofat{\edge}{\headvariableof{\edge}}$ is the controlled unitary
    \begin{align*}
        \qcaencodingofat{\acttensorof{\edge}}{\cavariableof{\edge}{\insymbol},\cavariableof{\edge}{\outsymbol},\headvariableof{\edge}}
        = \sum_{\headindexof{\edge}} \onehotmapofat{\headindexof{\edge}}{\headvariableof{\edge}} \otimes
        \yrotationofat{h\left(\frac{\acttensorofat{\edge}{\indexedheadvariableof{\edge}}}{
            \max_{\secheadindexof{\edge}}\acttensorofat{\edge}{\headvariableof{\edge}=\secheadindexof{\edge}}
        }\right)}{\cavariable{\insymbol},\cavariable{\outsymbol}} \, .
    \end{align*}
    The \activationCircuit{} of a tensor network $\{\acttensorofat{\edge}{\headvariableof{\edge}}\wcols\edgein\}$ is the contraction of the activation tensors to each tensor in the network, that is
    \begin{align*}
        \qcaencodingofat{\{\acttensorofat{\edge}{\headvariableof{\edge}}\wcols\edgein\}}{\cavariableof{\edges}{\insymbol},\cavariableof{\edges}{\outsymbol},\headvariableof{\nodes}}
        = \contractionof{
            \{\qcaencodingofat{\acttensorof{\edge}}{\cavariableof{\edge}{\insymbol},\cavariableof{\edge}{\outsymbol},\headvariableof{\edge}}\wcols\edgein\}
        }{\cavariableof{\edges}{\insymbol},\cavariableof{\edges}{\outsymbol},\headvariableof{\nodes}} \, .
    \end{align*}
\end{definition}

\begin{figure}[t]
    \begin{center}
        \input{tikz_pics/sampling/can_elementary_augmentation}
    \end{center}
    \caption{Ancilla augmentation of an elementary \ComputationActivationNetwork{} a) by replacing each leg vector of the elementary activation tensor by an directed ancilla augmentation.}
    \label{fig:canElementaryAugmentation}
\end{figure}

%% Elementary CAN Augmentation
Elementary \ComputationActivationNetworks{} can be augmented by ancilla augmentation of each leg vector in the activation tensor network (see \figref{fig:canElementaryAugmentation}).
We understand ancilla variables as variables in a Bayesian Network having single parents, corresponding with the ancilla augmentation of an elementary tensor.
For an example, see \figref{fig:toyAccounting}.

%% Hidden variables in activation graph
To a generic activation hypergraph $\graph$, we would do ancilla augmentation for the activation tensor network, where each hidden variable is treated as a distributed qubit.
To treat it as a distributed qubit, each hidden activation qubit is prepared in uniform state (i.e. a Hadamard gate acting on a ground state) and will be omitted in the measurement scheme (either not measured, or its measurement ignored).
This can be understood from a tensor network perspective, where marginalization means contracting, which is exactly how hidden activation variables are used.
For an example in the $\ttformat$ format, see \figref{fig:canTTAugmentation}.

\begin{figure}[t]
    \begin{center}
        \input{tikz_pics/sampling/can_TT_augmentation}
    \end{center}
    \caption{Ancilla augmentation of an $\ttformat$ \ComputationActivationNetwork{} a) by a Bayesian Network b).
    The hidden variables in the $\ttformat$ activation tensor are treated as distributed variables.
    Marginalization over the hidden variables and conditioning on the activation variables in state $1$, we reproduce the \ComputationActivationNetwork{}.}
    \label{fig:canTTAugmentation}
\end{figure}

\subsection{Q-samples of \CompActNets{}}


\begin{theorem}
    Given a \CompActNet{} with statistic $\sstat$ and activation tensor network $\acttensorat{\headvariables}$, the circuit composed of
    \begin{itemize}
        \item Hadamard transforms on the distributed qubits $\shortcatvariables$
        \item Computation circuit $\qcbencodingof{\sstat}$
        \item Activation circuit to $\qcaencodingof{\acttensorat{\headvariables}}$
    \end{itemize}
    prepares the q-sample of a distribution $\secprobat{\headvariables,\shortcatvariables,\avariables}$, which is conditioned on $\avariables=\onetuple{[\seldim]}$ the \CompActNet{}, that is
    \begin{align*}
        \secprobat{\headvariables,\shortcatvariables,\avariables=\onetuple{[\seldim]}}
        = \normalizationof{\bencodingofat{\sstat}{\headvariables,\shortcatvariables},\acttensorat{\headvariables}}{\headvariables,\shortcatvariables} \, .
    \end{align*}
\end{theorem}

\input{examples/sampling/toy_accounting_circuit}

% Bayesian network perspective
By the ancilla augmentation, a \CompActNet{} is turned into a Bayesian network, which has additional sparse structure in its computation network.
We can therefore understand the \computationCircuits{} and \activationCircuits{} as q-sample preparating circuits as in \cite{low_quantum_2014}.
The main difference is that we apply sparsity principles to further decompose the \computationCircuits{}.